\section{Foundations and related work}
\label{sec:foundations}

\subsection{Implicit surfaces and signed distance fields}

A signed distance field assigns to every point $\mathbf{x}\in\R^3$ the
signed Euclidean distance to the boundary $\partial\Omega$ of a solid
$\Omega$:
\begin{equation}
  f(\mathbf{x}) =
  \begin{cases}
    +\operatorname{dist}(\mathbf{x},\partial\Omega), & \mathbf{x}\notin\Omega,\\[2pt]
    -\operatorname{dist}(\mathbf{x},\partial\Omega), & \mathbf{x}\in\Omega.
  \end{cases}
  \label{eq:sdf}
\end{equation}
The surface is the zero set $f^{-1}(0)$.  The gradient of a true distance
field has unit length and points away from the solid, so the unit normal is
\begin{equation}
  \mathbf{n}(\mathbf{x}) = \frac{\nabla f(\mathbf{x})}{\|\nabla f(\mathbf{x})\|}.
  \label{eq:normal}
\end{equation}
The engine stores only eight corner samples per octree cell and
reconstructs $f$ inside the cell by trilinear interpolation.  This is a
deliberate approximation: it makes the field cheap to evaluate, cheap to
cache, and identical between the CPU reference and the GPU shader.

\subsection{CSG operations}

A CSG operation combines two distance fields.  The engine supports union,
subtraction, intersection, and a generic user operation.  For two fields
$f_1,f_2$ the common rules are
\begin{align}
  \text{union}:&\quad \min(f_1,f_2), \label{eq:union}\\
  \text{intersection}:&\quad \max(f_1,f_2), \label{eq:intersection}\\
  \text{subtraction}:&\quad \max(f_1,-f_2). \label{eq:subtraction}
\end{align}
The exact expressions are not distance fields in the strict sense after
subtraction (the result can be negative in the subtracted volume), but they
preserve the zero set and the sign, which is all the mesher needs.  The
engine also supports smooth minima and custom effects; those change the
field values but not the extraction algorithm.

\begin{figure}[htbp]
\centering
\begin{tikzpicture}[scale=1.0]
  \foreach \dx/\lab in {0/Union, 4.2/Intersection, 8.4/Subtraction, 12.6/Flip} {
    \begin{scope}[shift={(\dx,0)}]
      \draw[grid] (0,0) grid[step=0.5] (3,2.2);
      \fill[blue!12] (0.3,0.3) circle (0.75);
      \fill[blue!12] (1.8,1.1) circle (0.75);
      \ifdim\dx pt=0pt
        \draw[accent,thick] (0.3,0.3) circle (0.75);
        \draw[accent,thick] (1.8,1.1) circle (0.75);
      \fi
      \ifdim\dx pt=4.2pt
        \begin{scope}
          \clip (0.3,0.3) circle (0.75);
          \fill[accent2!25] (1.8,1.1) circle (0.75);
          \draw[accent!60,thick] (0.3,0.3) circle (0.75);
          \draw[accent2,thick] (1.8,1.1) circle (0.75);
        \end{scope}
      \fi
      \ifdim\dx pt=8.4pt
        \begin{scope}
          \clip (1.8,1.1) circle (0.75);
          \fill[white] (0.3,0.3) circle (0.75);
          \draw[accent!60,thick] (0.3,0.3) circle (0.75);
          \draw[accent2,thick] (1.8,1.1) circle (0.75);
        \end{scope}
      \fi
      \ifdim\dx pt=12.6pt
        \draw[accent,thick] (0.3,0.3) circle (0.75);
        \draw[accent2,thick,decorate,
          decoration={snake,amplitude=0.7mm,segment length=2mm}] (1.8,1.1) circle (0.75);
      \fi
      \node[below,font=\small] at (1.5,-0.15) {\lab};
    \end{scope}
  }
\end{tikzpicture}
\caption{Two-dimensional cross-sections of the CSG rules.  Blue is the
solid; a subtraction removes the second field from the first; the rightmost
panel illustrates a non-distance-preserving effect whose zero set and sign
are still well defined.}
\label{fig:csg}
\end{figure}

\subsection{Dual contouring and Surface Nets}

Marching Cubes places one vertex on each sign-changing cube edge and
connects them inside the cube.  Dual methods instead place \emph{one}
vertex per surface cell and connect the vertices of cells that share an
edge.  The resulting mesh is the dual of the cell complex.  Surface Nets
use an average of edge crossings; Dual Contouring solves a quadratic error
function (QEF) to place a vertex that fits the sampled Hermite data
(position and normal on each crossing edge).

The engine uses a QEF solve (with a mass-point regularizer) to place one
vertex per surface cell, then extracts triangles by scanning the twelve
edges of each scan cell.  The method is often called ``dual contouring'',
but the extraction rule is the Surface Nets ring rule: for each active
edge, collect the four cells around the edge, keep the surface cells,
order them around the edge, and triangulate the resulting three- or
four-gon.  The QEF is what makes the vertex position smooth; the ring rule
is what makes the connectivity exact and crack-free.

\subsection{Adaptive octrees and LoD}

A uniform dual grid at the finest resolution is wasteful: terrain near the
camera needs centimetres, while a distant ridge needs tens of metres.
An adaptive octree stores a cell only where the field crosses its cube and
subdivides cells near the surface.  Each cell has a ``level''; a cell is
allowed to be a leaf (no children) when its cube is small enough or when
its eight children can be collapsed without exceeding an error threshold.

Adaptive dual meshing introduces two problems:

\begin{enumerate}
\item \textbf{The dual vertex of a cell can be far from its cube.}  A QEF
  vertex is clamped inside its cell, but the four cells around an edge can
  have different sizes.
\item \textbf{Two neighbouring chunks can choose different levels.}  A
  coarse cell on one side of a boundary and a fine cell on the other do
  not share the same edge vertices, so a naive fan leaves a crack.
\end{enumerate}

The engine solves both with the same idea: for every active edge, sample
the four quadrant cells at the edge midpoint and recursively collect
\emph{all} cell-boundary coordinates along the edge.  The edge is split at
those coordinates, and each sub-segment is emitted by its owner cell.  A
neighbouring chunk resolves the same edge with the same owner and the same
sub-segment list, so the two sides agree bit-for-bit.

\subsection{Why two processors?}

The split used here is not a generic ``CPU authoring, GPU rendering''
division.  It is driven by the representation:

\begin{itemize}
\item \textbf{The octree is a mutable, shared, pointer-rich structure.}
  It is edited by brushes, resampled by shape operations, versioned,
  simplified, and guarded by a read/write lock.  This is natural on the CPU
  and hostile to a GPU.
\item \textbf{The triangle extraction is data-parallel and pure.}  Given a
  cell's eight SDF samples and the four quadrant cells of an edge, the
  output is a deterministic function.  This is ideal for a compute shader.
\item \textbf{The renderer wants device-local buffers.}  A CPU mesher
  produces a host array, which must be copied; a GPU mesher can write the
  renderer's buffers directly.
\end{itemize}

The cost is a flattening step.  The host must traverse the octree
neighbourhood and produce a table small enough to upload per chunk.  The
flattening is exact: it replays the same root-to-cell descent and the same
child-index arithmetic the shader will use, so the GPU sees a
\emph{serialisable snapshot} of the tree, not a lossy approximation.

\subsection{Position in the literature}

The meshing rule descends from Gibson's Surface Nets and Ju et al.'s Dual
Contouring.  The adaptive stitching rule is related to the dual-cell
``edge split'' methods used in transvoxel-like terrain engines, but it is
derived here directly from the CPU walker's break collection rather than
from a fixed stitching table.  The GPU port follows the same arithmetic
order as the CPU reference (including the QEF's Cramer solve and the
guarded trilinear lerp) so the two can be compared and the offline exporter
can reproduce the runtime mesh.
